% ch1_d §1.6 (书页 23–26 / p037–p040)
\section{边界条件：态的计数}

以上这一切都假定了完全的平移对称性，而这将要求晶格为无限。这在数学上相当麻烦，因为那样一来，原子和有待处理的态都会有无限多个。我们唯一的办法，是先处理一个原子数目有限的系统，然后令这个数目趋于无穷而取极限。我们也许会像对待真实表面那样引入边界，例如要求波函数在边界上必须为零。但这样做同样不便，因为此时系统的严格定态将不得不计及来自边界的所有反射，结果都会成为驻波。传导电子通常在极短的距离内就会被杂质或热振动非相干地散射开；要用这类函数去表示传导电子的态，在数学上笨拙，而且在本质上是不物理的。

%FIG{14}: 简约区中自由电子的能量。

有一个数学上的手法，能圆满地解决态的计数问题，却又不会引入任何来自边界的直接物理效应。这就是采用循环（cyclic）边界条件，即 Born–von Kármán 边界条件。

在一维情形，设这个“晶体”含有 $L$ 个晶胞，它们首尾相接连成一个环。这一做法对电子波函数的后果，可以用如下条件来表达：
\begin{equation*}
\psi(x+La)=\psi(x),
\tag{1.71}
\end{equation*}
这就保证了在接合点处没有不连续。

但我们的一维 Bloch 条件给出
\begin{equation*}
\psi_k(x+La)=e^{ikLa}\psi_k(x),
\tag{1.72}
\end{equation*}
于是必须有
\begin{equation*}
e^{ikLa}=1,
\tag{1.73}
\end{equation*}
即
\begin{equation*}
k=\frac{2\pi m}{La},
\tag{1.74}
\end{equation*}
其中 $m$ 为整数。

在一维的简约区中，我们有 $-\pi/a<k\leqslant\pi/a$，因此落在范围
\begin{equation*}
-\frac{1}{2}L<m\leqslant\frac{1}{2}L
\tag{1.75}
\end{equation*}
之内的那些整数，将给出所有本质上不同的简约波数值。这样的值恰好有 $L$ 个（这里当然要分别处理 $L$ 为奇数和 $L$ 为偶数的不同情形——一个琐碎的细节），并且相邻两值相隔 $(1/L)(2\pi/a)$。由于假定 $L$ 非常大，这个分布实际上是连续的，并且在倒格子空间中具有恒定的密度。

在三维情形，我们把这一论证加以推广，宣称宏观系统在三个维度上都是循环的。我们把它取为沿晶格三个基矢方向尺寸分别为 $L_1\mathbf{a}_1$、$L_2\mathbf{a}_2$、$L_3\mathbf{a}_3$ 的晶体，并写下
\begin{equation*}
\psi(\mathbf{r}+L_1\mathbf{a}_1)=\psi(\mathbf{r}),\quad
\psi(\mathbf{r}+L_2\mathbf{a}_2)=\psi(\mathbf{r}),\quad
\psi(\mathbf{r}+L_3\mathbf{a}_3)=\psi(\mathbf{r}),
\tag{1.76}
\end{equation*}
正如式 (1.71) 一样。

对于波矢为 $\mathbf{k}$ 的 Bloch 态，这些条件意味着
\begin{equation*}
e^{i\mathbf{k}\cdot(L_1\mathbf{a}_1)}
=e^{i\mathbf{k}\cdot(L_2\mathbf{a}_2)}
=e^{i\mathbf{k}\cdot(L_3\mathbf{a}_3)}=1,
\tag{1.77}
\end{equation*}
而这只有当 $\mathbf{k}$ 取如下形式时方能实现：
\begin{align*}
\mathbf{k} &= k_1\mathbf{b}_1+k_2\mathbf{b}_2+k_3\mathbf{b}_3\\
&= \frac{2\pi m_1}{L_1}\mathbf{b}_1+\frac{2\pi m_2}{L_2}\mathbf{b}_2+\frac{2\pi m_3}{L_3}\mathbf{b}_3,
\tag{1.78}
\end{align*}
其中 $m_1$、$m_2$、$m_3$ 为整数，而矢量 $\mathbf{b}_1$、$\mathbf{b}_2$、$\mathbf{b}_3$ 当然就是 $(\mathbf{a}_1,\mathbf{a}_2,\mathbf{a}_3)$ 的倒格子三矢组（reciprocal triad），如式 (1.19) 中那样。

但这一结果显然可与倒格矢的定义相类比：
\begin{equation*}
\mathbf{g}=n_1\cdot 2\pi\mathbf{b}_1+n_2\cdot 2\pi\mathbf{b}_2+n_3\cdot 2\pi\mathbf{b}_3,
\tag{1.79}
\end{equation*}
其中 $n_1$、$n_2$、$n_3$ 为整数。把倒格子晶胞的生成矢量沿 $\mathbf{b}_1$ 方向分成 $L_1$ 份、沿 $\mathbf{b}_2$ 方向分成 $L_2$ 份、沿 $\mathbf{b}_3$ 方向分成 $L_3$ 份，便得到式 (1.78) 中 $\mathbf{k}$ 的允许值（allowed values）。于是，宛如生出一阵细密的疹点，点子均匀地布满倒格子空间，如图 15 所示。

要计算这些点的密度，只需注意到：让整数 $m_1$、$m_2$、$m_3$ 取遍
\begin{equation*}
0\leqslant m_1<L_1,\quad
0\leqslant m_2<L_2,\quad
0\leqslant m_3<L_3.
\tag{1.80}
\end{equation*}
即可铺满倒格子的一个晶胞。

但这并不是我们自然会选来作为 $\mathbf{k}$ 矢量基本区的晶胞。或许改取如下范围更好：
\begin{equation*}
-\frac{1}{2}L_1<m_1<\frac{1}{2}L_1,\quad
-\frac{1}{2}L_2<m_2<\frac{1}{2}L_2,\quad
-\frac{1}{2}L_3<m_3<\frac{1}{2}L_3,
\tag{1.81}
\end{equation*}
以与式 (1.75) 类比。这将给出一个以原点为中心的平行六面体晶胞。然而，晶胞的最佳选择是倒格子的 Wigner–Seitz 原胞，也就是我们的老朋友——布里渊区。

%FIG{15}: 布里渊区与倒格子的平行六面体晶胞覆盖相同的面积，且恰含 N 个“允许的 k 矢量”。

由于布里渊区的体积与平行六面体晶胞相同，它必定恰好含有同样数目的 $\mathbf{k}$ 的“允许值”。由式 (1.80) 或式 (1.81)，这个数目显然就是 $L_1\times L_2\times L_3=N$，而这正是我们整块宏观晶体中晶胞的数目。这是一个极其重要的定理：\emph{一块晶体中有多少个晶胞，一个布里渊区中就恰有多少个允许的波矢。}

我们已经证明过（式 (1.32)），一个区的体积为 $8\pi^3/v_c$。若体积为 $V$ 的晶体中有 $N$ 个晶胞，则
\begin{equation*}
v_c=V/N.
\tag{1.82}
\end{equation*}
于是，每个允许的 $\mathbf{k}$ 矢量在 $\mathbf{k}$ 空间中占据的体积便是
\begin{equation*}
\frac{1}{N}\,\frac{8\pi^3}{v_c}=\frac{8\pi^3}{V}.
\tag{1.83}
\end{equation*}
因此，\emph{倒格子空间单位体积内有 $V/8\pi^3$ 个允许的 $\mathbf{k}$ 矢量。}

实际上 $N$ 非常大，以致这一分布被当作连续分布来处理。我们常把对 $\mathbf{k}$ 矢量的求和写成一个积分
\begin{equation*}
\sum_{\mathbf{k}}\rightarrow\int d\mathbf{k}\equiv\frac{V}{8\pi^3}\iiint d^3k,
\tag{1.84}
\end{equation*}
这里用单个积分作为求和之极限的简明记号。然而，重要的是要记住：当我们真正要在 $\mathbf{k}$ 空间中做具体求积时，必须计入权重因子 $V/8\pi^3$。为简单起见，我们通常将假定 $V=1$，这时 $N$ 就是单位体积晶体中的晶胞数，而 $1/N$ 就是一个晶胞的体积。

结果 (1.83) 实际上不依赖于任何所假定的区结构；它在辐射理论中以及在自由电子情形下早已为人熟知。但就许多目的而言，布里渊区恰好含有 $N$ 个允许点这一事实更为有用。它表明这个区在本质上是不变的，仅依赖于晶体结构；增大整块晶体的尺寸，只会使 k 空间中的态密度（density of states）增大。

不过，必须承认，Born–von Kármán 条件在物理上是无法实现的。在二维情形，画在环面（torus）表面上的晶胞网络沿两个方向都是循环的；但对三维晶格来说，无论怎样施展拓扑上的扭曲，都无法使三个条件同时满足。尽管如此，它作为一种数学技巧用起来极为出色，而且可以用一些繁冗的数学来为它提供证明。事实是：在大量子数的渐近极限下，态密度对边界条件的具体形式非常不敏感。只要我们并不打算讨论边界效应本身，循环边界条件就远远是最方便的选择。
